\documentclass[12pt]{article}
% \usepackage[utf-8]{inputenc}
\usepackage[margin=1in]{geometry}
\usepackage{listings}
\usepackage{xcolor}
\usepackage{graphicx}
\usepackage{hyperref}
\usepackage{amsmath}
\usepackage{amssymb}
\usepackage{booktabs}

% Configure listings style
\lstset{
    language=Python,
    basicstyle=\ttfamily\small,
    keywordstyle=\color{blue},
    commentstyle=\color{gray},
    stringstyle=\color{red},
    breaklines=true,
    breakatwhitespace=true,
    showstringspaces=false,
    tabsize=4,
    backgroundcolor=\color{gray!10},
    frame=single,
    linewidth=0.95\textwidth
}

\title{Homework: Multithreaded Server with a Thread Pool}
\author{Haotang Li}
\date{}

\begin{document}

\maketitle

\section{Implementation}

\subsection{Thread Pool Architecture}

\begin{lstlisting}
def start(self) -> None:
    # Create fixed number of worker threads
    for i in range(self.workers):
        t = threading.Thread(target=self._worker_loop, args=(i,), daemon=True)
        self._threads.append(t)
        t.start()
\end{lstlisting}

\textbf{Worker Loop:}
\begin{lstlisting}
def _worker_loop(self, worker_id: int) -> None:
    while not self._stop.is_set():
        task = self.task_q.get()
        if task.conn is None:  # Sentinel for shutdown
            self.task_q.task_done()
            break
        
        resp = parse_and_execute(task.line)
        # Send response with per-connection lock
        with conn_lock:
            task.conn.sendall(resp_line.encode("utf-8"))
        
        self.task_q.task_done()
\end{lstlisting}

\textbf{Design Rationale:}

\begin{itemize}
    \item \textbf{Fixed pool size}: Avoids thread creation overhead and OS resource exhaustion
    \item \textbf{Daemon threads}: Automatically terminated when main thread exits
    \item \textbf{Worker ID}: Enables per-worker debugging and monitoring
    \item \textbf{Sentinel pattern}: Clean shutdown by injecting \texttt{None} tasks to wake blocked workers
\end{itemize}

\pagebreak

\subsection{Bounded Queue}

\begin{lstlisting}
self.task_q: "queue.Queue[Task]" = queue.Queue(maxsize=queue_size)
\end{lstlisting}

\textbf{Command-line Control:}
\begin{verbatim}
python server.py --queue 500  # Set queue size to 500
\end{verbatim}

\textbf{Queue Characteristics:}

\begin{itemize}
    \item \textbf{Thread-safe}: Python's \texttt{queue.Queue} handles internal locking
    \item \textbf{Bounded}: Maximum capacity prevents memory exhaustion under high load
    \item \textbf{Blocking operations}: \texttt{put()} and \texttt{get()} support timeout-based backpressure
\end{itemize}

\textbf{Memory Safety:}
\begin{itemize}
    \item At queue size 500 with \textasciitilde200 bytes per task: \textasciitilde100KB memory footprint
    \item Prevents OOM conditions during traffic spikes
    \item Configurable limit allows tuning for specific workloads
\end{itemize}

\subsection{Backpressure Mechanism}

\textbf{Implementation - Two Modes:}

\subsubsection{Mode 1: Block with 1-second timeout (default)}

\begin{lstlisting}
try:
    self.task_q.put(task, block=True, timeout=1.0)
except queue.Full:
    conn.sendall(b"ERR server busy\n")
\end{lstlisting}

\subsubsection{Mode 2: Immediate rejection}

\begin{verbatim}
python server.py --reject-when-full
\end{verbatim}

\begin{lstlisting}
try:
    self.task_q.put_nowait(task)
except queue.Full:
    conn.sendall(b"ERR server busy\n")
\end{lstlisting}

\textbf{Thread Safety:}
\begin{itemize}
    \item Error responses use per-connection locks to avoid race conditions
    \item Prevents interleaved writes when multiple threads reject requests
\end{itemize}

\textbf{Behavior Under Load:}
\begin{itemize}
    \item \textbf{Mode 1}: Provides brief buffering window for transient spikes
    \item \textbf{Mode 2}: Immediate feedback, better for latency-sensitive clients
    \item Both modes prevent queue overflow and provide clear error signals
\end{itemize}

\subsection{Concurrency Handling}

\subsubsection{A. Multiple Concurrent Clients}

\textbf{Accept Loop:}
\begin{lstlisting}
def _accept_loop(self) -> None:
    server_sock = socket.socket(socket.AF_INET, socket.SOCK_STREAM)
    server_sock.setsockopt(socket.SOL_SOCKET, socket.SO_REUSEADDR, 1)
    server_sock.bind((self.host, self.port))
    server_sock.listen()
    
    while not self._stop.is_set():
        conn, addr = server_sock.accept()
        t = threading.Thread(target=handle_client, args=(conn, addr), daemon=True)
        t.start()
\end{lstlisting}

\textbf{Per-client threads}: Each connection handled in dedicated thread for isolation

\subsubsection{B. Connection Lifecycle Management}

\textbf{Client Handler:}
\begin{lstlisting}
def handle_client(conn: socket.socket, addr: Tuple[str, int]) -> None:
    try:
        conn.setsockopt(socket.IPPROTO_TCP, socket.TCP_NODELAY, 1)
        buf = b""
        while not self._stop.is_set():
            data = conn.recv(4096)
            if not data:  # Client closed connection
                break
            # Process buffered lines...
    finally:
        try:
            conn.close()
        except OSError:
            pass
        # Clean up connection lock
        with self._conn_locks_lock:
            self._conn_locks.pop(conn, None)
\end{lstlisting}

\textbf{Key Features:}
\begin{itemize}
    \item \textbf{TCP\_NODELAY}: Reduces latency by disabling Nagle's algorithm
    \item \textbf{Graceful shutdown}: \texttt{finally} ensures socket cleanup
    \item \textbf{Memory leak prevention}: Connection locks removed on close
\end{itemize}

\subsubsection{C. Thread-Safe Response Transmission}

\textbf{Problem}: Client pipelining causes multiple workers to write to same connection concurrently.

\textbf{Solution - Per-Connection Locks:}
\begin{lstlisting}
# In __init__
self._conn_locks: dict[socket.socket, threading.Lock] = {}
self._conn_locks_lock = threading.Lock()

# In _worker_loop
with self._conn_locks_lock:
    if task.conn not in self._conn_locks:
        self._conn_locks[task.conn] = threading.Lock()
    conn_lock = self._conn_locks[task.conn]

with conn_lock:
    task.conn.sendall(resp_line.encode("utf-8"))
\end{lstlisting}

\textbf{Benefits:}
\begin{itemize}
    \item \textbf{Fine-grained locking}: Different connections can write concurrently
    \item \textbf{Serialization per connection}: Prevents response interleaving
    \item \textbf{Dynamic lock creation}: Locks allocated only for active connections
\end{itemize}

\section{Testing}

\subsection{Test Environment}

\begin{itemize}
    \item \textbf{Platform}: macOS
    \item \textbf{Python}: 3.12.12
    \item \textbf{Network}: localhost (127.0.0.1)
\end{itemize}

\subsection{Test Case 1: Basic Functionality}

\textbf{Objective}: Verify correct request/response handling

\textbf{Configuration:}
\begin{verbatim}
python server.py --host 127.0.0.1 --port 9000 --workers 8 --queue 500

python3 client.py --host 127.0.0.1 --port 9000 --clients 5 --requests 200 --mix balanced
\end{verbatim}

\textbf{Result:}

Server:
\begin{verbatim}
[stats] processed=450 avg_latency_ms=5349.49 qlen=0 rps=18.71
\end{verbatim}

Client:
\begin{verbatim}
[INFO] [MainThread] root: SUMMARY sent=450 recv=200 ok=200 err=0 elapsed=11.72s ok_rps=17.07
\end{verbatim}

\textbf{Analysis}: The max-inflight of client in default is 50, and the timeout is 5 sec, exceeding the concurrency capacity of server.
\begin{verbatim}
[stats] processed=425 avg_latency_ms=5215.63 qlen=17 rps=26.51
\end{verbatim}

The log shows the maximum capacity of server is about 27 rps, and the avg\_latency\_ms is about 5.5 sec.

\subsection{Test Case 2: High Concurrency of workers}

\textbf{Objective}: Test scalability with multiple concurrent workers of server.

\textbf{Configuration:}
\begin{verbatim}
python server.py --host 127.0.0.1 --port 9000 --workers 32 --queue 500

python3 client.py --host 127.0.0.1 --port 9000 --clients 5 --requests 200 --mix balanced
\end{verbatim}

\textbf{Result:}

Server:
\begin{verbatim}
[stats] processed=853 avg_latency_ms=1110.70 qlen=115 rps=106.41
[stats] processed=1000 avg_latency_ms=1154.88 qlen=0 rps=99.83
\end{verbatim}

Client:
\begin{verbatim}
[INFO] [MainThread] root: SUMMARY sent=1000 recv=1000 ok=1000 err=0 elapsed=8.56s ok_rps=116.77
\end{verbatim}

\textbf{Analysis}:
\begin{itemize}
    \item All 1000 requests processed successfully
    \item Peak RPS: 106.41 req/s
    \item Queue never overflowed
    \item No client timeouts with reasonable timeout settings
\end{itemize}

\subsection{Test Case 3: Backpressure - Immediate Reject Mode}

\textbf{Objective}: Test immediate rejection when queue full

\textbf{Configuration:}
\begin{verbatim}
python server.py --host 127.0.0.1 --port 9000 --workers 8 --queue 500 --reject-when-full

python3 client.py --host 127.0.0.1 --port 9000 --clients 5 --requests 200 --mix balanced
\end{verbatim}

\textbf{Result:}

Server:
\begin{verbatim}
[stats] processed=375 avg_latency_ms=5563.65 qlen=17 rps=23.41
[stats] processed=400 avg_latency_ms=5664.21 qlen=0 rps=22.19
\end{verbatim}

Client:
\begin{verbatim}
[INFO] [MainThread] root: SUMMARY sent=400 recv=150 ok=150 err=0 elapsed=11.71s ok_rps=12.81
\end{verbatim}

\textbf{Analysis}:
\begin{itemize}
    \item Confirms immediate rejection is more aggressive
    \item Appropriate for latency-sensitive applications
\end{itemize}

\subsection{Scaling Analysis}

\textbf{Thread Pool Size Impact:}
\begin{verbatim}
Workers=2:  ~5.5 req/s
Workers=4:  ~13.7 req/s
Workers=8:  ~27 req/s
Workers=16: ~52.5 req/s
Workers=32: ~102.4 req/s
\end{verbatim}

\textbf{Observation}: Near-linear scaling for I/O-bound workloads up to 32 workers.

\end{document}
